\section{Results}
\label{sec:results}

Every number in this section is harvested mechanically from the submitted output
files. The convergence status quoted for each case is the status the solver itself
wrote to \texttt{run\_status.json}; no case is relabelled here. Where a value
appears as \pending, that run had not completed when this document was last
compiled, and the entry is deliberately left visibly empty rather than filled
with an estimate.

\subsection{Summary Tables}
\label{sec:summary}

\Cref{tab:runstatus} is the run-status table and \cref{tab:forces} the force
table. In \cref{tab:forces} the final column reports the measured stationarity of
the drag history: the change in $C_D$ across the trailing tenth of the recorded
history, classified against the criterion of \cref{eq:stationarity}. A case marked
\textbf{drifting} has not reached a steady state, regardless of how many residual
orders it achieved, and its coefficients should be read as a snapshot rather than
a converged value.

\begin{table}[htbp]
  \centering
  \caption{Run status for all eight required cases, transcribed from each
  \texttt{run\_status.json}. Physical time is zero for steady cases by
  construction. Wall time is for the rank count shown, under the 4-CPU container
  quota described in \cref{sec:quota}.}
  \label{tab:runstatus}
  \footnotesize
  \begin{tabular}{lrrrrrl}
\toprule
Case & Ranks & Steps & $t_{\mathrm{phys}}$ & Orders & Wall (s) & Status \\
\midrule
NACA $M{=}0.15$ inv. & \num{4} & \num{2225} & \num{0} & \num{4.71} & \num{29.64} & converged \\
NACA $M{=}0.80$ inv. & \num{4} & \num{3496} & \num{0} & \num{2.89} & \num{67.19} & converged \\
NACA $M{=}2.0$ inv. & \num{4} & \num{4906} & \num{0} & \num{2.35} & \num{48.66} & converged \\
NACA $M{=}0.15$ lam. & \num{4} & \num{3839} & \num{0} & \num{5.47} & \num{53.3} & converged \\
NACA $M{=}0.80$ lam. & \num{4} & \num{4639} & \num{0} & \num{4} & \num{55.29} & converged \\
NACA $M{=}2.0$ lam. & \num{4} & \num{3646} & \num{0} & \num{3.88} & \num{34.55} & converged \\
Cyl.\ $Re{=}20$ & \num{4} & \num{2450} & \num{0} & \num{5.86} & \num{21.83} & converged \\
Cyl.\ $Re{=}200$ & \pending &  &  &  &  &  \\
\bottomrule
\end{tabular}

\end{table}

\begin{table}[htbp]
  \centering
  \caption{Force coefficients from the final row of each \texttt{forces.csv}, with
  the pressure/viscous drag split. The final column records how the run established
  stationarity, taken from the solver's own termination record rather than
  re-derived, and distinguishes three outcomes: a \emph{fixed point} reached within
  the span tolerance; a \emph{limit cycle} whose mean is converged but whose state
  keeps oscillating; and an \emph{asymptotic} approach, still monotone at
  termination but with geometrically decaying increments, for which the quoted
  $\pm$ is the solver's estimate of the $C_D$ movement still remaining and therefore
  the precision of the coefficient (see \cref{sec:digits}). For limit cycles the
  oscillation
  amplitude is given, as a percentage of $C_D$ where that is meaningful and in
  absolute terms for the near-zero-drag case. \textbf{For the limit-cycle rows the
  quoted $C_D$ is a single-instant snapshot and its trailing digits are not
  significant}; the meaningful statement for those cases is the mean plus amplitude,
  given in \cref{eq:m200drag} for $\Minf=2.0$ inviscid and in
  \cref{sec:naca} for the others. The $\Reinf=200$ transient is periodic by
  construction, so the steady stationarity tests do not apply to it: its drag is
  \emph{supposed} to oscillate, and its convergence is established by the periodicity
  measurements of \cref{tab:saturation} instead. Its row is a single-instant snapshot of
  a periodic signal; the mean and amplitude given in \cref{sec:cyl200} are the
  meaningful quantities.}
  \label{tab:forces}
  \footnotesize
  \setlength{\tabcolsep}{4pt}
  \begin{tabular}{lrrrrrl}
\toprule
Case & $C_D$ & $C_L$ & $C_{m,z}$ & $C_{D,p}$ & $C_{D,v}$ & Stationarity \\
\midrule
NACA $M{=}0.15$ inv. & \num{0.0009604} & \num{0.000156} & \num{6.82e-05} & \num{0.0009604} & \num{0} & cycle, 12.9\% p--p \\
NACA $M{=}0.80$ inv. & \pending &  &  &  &  &  \\
NACA $M{=}2.0$ inv. & \pending &  &  &  &  &  \\
NACA $M{=}0.15$ lam. & \pending &  &  &  &  &  \\
NACA $M{=}0.80$ lam. & \pending &  &  &  &  &  \\
NACA $M{=}2.0$ lam. & \pending &  &  &  &  &  \\
Cyl.\ $Re{=}20$ & \num{2.018} & \num{0.000517} & \num{-2.07e-06} & \num{1.221} & \num{0.7975} & asymptotic, $\pm$\num{9.89e-04} \\
Cyl.\ $Re{=}200$ & \pending &  &  &  &  &  \\
\bottomrule
\end{tabular}

\end{table}

\subsection{How many digits of each coefficient are earned}
\label{sec:digits}

An adversarial review of the completed runs raised a question that deserves a direct
answer rather than a defence: several cases were stopped while their drag history
was still moving in one direction, so how much of each quoted coefficient is real?
The answer turned out to require a distinction the solver was not making, and it now
makes it and reports the result, so this subsection describes a criterion in the code
rather than a post-hoc analysis of it.

The question cannot be answered by the convergence status alone, because
\emph{monotone motion and non-convergence are not the same thing}. What separates
them is how the increments behave. Take the sub-block means of $C_D$ over the
trailing window and difference them:

\begin{itemize}
  \item If the successive decrements shrink \textbf{geometrically}, with a roughly
  constant ratio $r<1$, the iteration is in its asymptotic regime. The distance
  still to be covered is the sum of the remaining geometric series,
  $\Delta_{\infty}\approx d_{\text{last}}\,r/(1-r)$, which is bounded, small, and
  computable. The run is converged, and the remaining error is quantifiable.
  \item If the decrements are roughly \textbf{constant}, the solution is marching at
  a steady rate and there is no asymptote in sight. That is an unfinished transient,
  and no extrapolation of it is defensible.
\end{itemize}

This distinction is what the review's monotone-trend observation actually detects,
and it separates the cases cleanly. For the cylinder at $\Reinf=20$ the sub-block
decrements over the certifying window run
\begin{center}
  $-4.42,\ -3.63,\ -2.94,\ -2.22,\ -1.64,\ -1.18,\ -0.82,\ -0.55,\ -0.36$
  \quad (units of $10^{-3}$),
\end{center}
with successive ratios $0.82, 0.81, 0.76, 0.74, 0.72, 0.70, 0.67, 0.65$. That is a
clean geometric decay, with a mean ratio of $0.733$.

The solver now performs this estimate itself and states the resulting precision in its
own termination note, so the figure quoted here is the solver's rather than an external
reconstruction of it: the geometric tail gives $\num{9.89e-4}$ of $C_D$ movement
remaining, or $0.049\,\%$, for an asymptote near $2.0174$. \textbf{Three digits of this
coefficient are earned and the fourth is not}, and it is quoted that way in
\cref{sec:cyl20}.

That figure has two independent corroborations, which is why it is quoted with
confidence. First, the extrapolation script written for this report, working only from
the submitted \texttt{forces.csv} and knowing nothing of the solver's internals,
independently gives an asymptote of $2.0177$ bracketed between $2.0176$ and $2.0178$
across the range of observed ratios --- agreeing with the solver's own estimate to
$\num{3e-4}$. Second, the sensitivity study of \cref{sec:sensitivity} ran this case
further, to $6.73$ residual orders, and obtained $C_D=\num{2.018119}$, which sits
between the reported value and the extrapolated asymptote, exactly where a longer run
on a decaying tail is required to land. Two runs of different length and two
independent extrapolations agree on the first three digits and disagree in the fourth,
which is a more useful statement of the uncertainty than any single run's trailing
digits.

This asymptote is quotable as a number because the extrapolation is well conditioned
here. By \cref{eq:tailsum} the fitted ratio's error is amplified by $r/(1-r)$, which at
$r=0.733$ is only $2.7$, so a percent-level error in the ratio moves the asymptote by a
few percent of an already small remaining distance. \Cref{sec:conditioning} shows the same
formula at $r=0.994$, where the amplification is $165$ and no asymptote may be quoted.
The distinction is quantitative rather than a matter of judgement, and it is the test
applied throughout this report before any extrapolated figure is given.

By contrast, the case the review identified as the worst instance,
\texttt{naca0012\_m200\_laminar\_re5000}, fails this test decisively: its decrement
ratios are $1.11, 1.27, 0.86, 0.85, 1.04, 1.03, 1.00, 0.99$ with a mean of $1.019$, so
the increments are not decaying at all and the tail sum diverges. There is no asymptote
to extrapolate to, and none is offered. Combined with the component cancellation of
\cref{sec:cancellation}, that case is reported as not converged.

The contrast between mean ratios of $0.733$ and $1.019$ is what makes the discriminator
usable rather than merely descriptive: the two cases are separated by a wide margin on a
single measured quantity, not by a judgement call. It is worth being explicit that this
vindicates part of the review's finding and corrects part of it. The review was right
that a monotone certifying window is a warning sign and right that
\texttt{naca0012\_m200\_laminar\_re5000} was not converged. It was wrong to conclude
that all monotone cases were unconverged: a monotone approach with geometrically
decaying increments is the normal, correct behaviour of a converging iteration, and a
criterion that rejected it would reject good runs along with bad ones.

The general policy adopted for this section follows from the above: coefficients are
quoted to the precision the convergence evidence supports, an extrapolated asymptote
and its bracket are given where the tail is demonstrably geometric, and the trailing
digits produced by the solver are not reproduced merely because they exist. The
script that performs the extrapolation is \texttt{report/extrapolate\_tail.py} and
the window diagnostics are \texttt{report/check\_window.py}, both of which read only
the submitted \texttt{forces.csv}.

\paragraph{Three different things limit precision, and each is quantified.} It is worth
collecting them, because they are independent and a reader should know which applies where:

\begin{itemize}
  \item \textbf{Limit-cycle amplitude.} For $\Minf=2.0$ inviscid the solution is a bounded
  oscillation, so no number of iterations produces more digits: the honest report is the mean
  plus the amplitude, $C_D=0.0880\pm0.0007$ (\cref{sec:limitcycle}).
  \item \textbf{Conditioning of the tail estimate.} For a monotone asymptotic approach the
  remaining movement is computable, but only where the decay ratio is away from unity. At
  $r=0.733$ the cylinder's asymptote is quotable; at $r=0.994$ nothing is
  (\cref{sec:conditioning}).
  \item \textbf{Partition sensitivity.} For the one case needing tens of thousands of
  iterations, two rank counts give answers $\num{5.0e-4}$ apart, far more than either run's
  own tail estimate, so its fourth digit is partition-dependent rather than converged
  (\cref{sec:conditioning}).
\end{itemize}

None of the three is a defect; each is a property of the problem or the algorithm, measured
rather than assumed. Their value is that they say where the reported digits come from, and
by implication that the coefficients not carrying such a caveat are limited only by the
convergence level stated beside them.

\subsection{NACA0012 Cases}
\label{sec:naca}

All six airfoil cases use the same mesh at zero angle of attack, so the exact
solution is symmetric about $y=0$ and the exact lift and moment are zero. The
computed $C_L$ and $C_{m,z}$ in \cref{tab:forces} are therefore measures of the
discretization's ability to preserve symmetry on a mesh that is not itself
symmetric in its cell connectivity.

One caveat must be attached to that statement immediately, because a near-zero
integrated $C_L$ is a weaker test than it appears. Lift is an \emph{integral} of the
surface pressure difference, so equal and opposite local asymmetries cancel in it. A
small $C_L$ is therefore necessary but not sufficient for a symmetric solution, and
the pointwise surface distribution is the stronger check. \Cref{sec:m200asym}
reports a case where the integrated lift is small while the local surface
distribution is badly asymmetric, so $C_L$ alone should not be read as evidence of
symmetry preservation for the supersonic case.

\subsubsection{Subsonic inviscid, $\Minf=0.15$}

This case is the cleanest available measure of pure discretization error. At
$\Minf=0.15$ with no shock and no viscosity, d'Alembert's paradox applies and the
exact drag is identically zero. Any computed $C_D$ is therefore entirely numerical
--- artificial dissipation from the Riemann solver plus reconstruction error --- so
its magnitude, not its sign, is the meaningful quantity. The solver reports
$C_D=\cnsCdNacaSubInv$ with $C_L=\cnsClNacaSubInv$, over
$\cnsOrdersNacaSubInv$ orders of residual reduction in $\cnsStepsNacaSubInv$
steps.

\Cref{fig:naca015inv-hist} shows the residual and force histories and
\cref{fig:naca015inv-field} the Mach and pressure fields. The $C_p$ distribution
in \cref{fig:naca015inv-cp} should be symmetric between upper and lower surfaces
and should approach the stagnation value at the leading edge. The surface row nearest
the nose reports $C_p=\cnsNoseCpNacaSubInv$, which \emph{exceeds} the familiar value of
$1$ --- and that is correct, because $C_p=1$ is the \emph{incompressible} stagnation
value. For compressible flow the isentropic stagnation pressure coefficient is
\begin{equation}
  C_{p0}=\frac{p_0/p_\infty-1}{\tfrac12\gamma \Minf^2}
  =\frac{\left(1+\tfrac{\gamma-1}{2}\Minf^2\right)^{\gamma/(\gamma-1)}-1}
         {\tfrac12\gamma \Minf^2},
  \label{eq:stagcp}
\end{equation}
which exceeds unity and tends to $1$ as $\Minf\to0$. It evaluates to
$\cnsStagCpNacaSubInv$ at $\Minf=0.15$ and $\cnsStagCpNacaTraInv$ at $\Minf=0.80$.
\Cref{eq:stagcp} is the right reference only while the stagnation streamline is isentropic,
which is true subsonically but \emph{not} at $\Minf=2.0$: there the flow crosses the bow
shock first, so the attainable ceiling is the lower post-shock pitot value of
\cref{eq:pitot}, $1.6573$, rather than the isentropic $\cnsStagCpNacaSupInv$. That case is
treated separately in \cref{sec:m200asym}.

The size of that distinction is worth one more sentence, because choosing the wrong
reference at $\Minf=2.0$ would not merely be imprecise. The isentropic value
$\cnsStagCpNacaSupInv$ overstates the attainable ceiling by $47.1\,\%$ relative to the
post-shock $1.6573$, and on the supersonic inviscid case $9$ of $404$ wall faces exceed the
correct bound while \textbf{zero} exceed the isentropic one. Quoting the isentropic figure
would therefore have converted a genuine, diagnostic bound violation into an apparently
clean result, and removed the evidence that \cref{sec:m200asym} rests on. A reference value
is not a presentational detail when it is the thing a bound is measured against.

Measured against the correct reference, the computed nose value is a deficit of
$\cnsNoseCpDeficitPctNacaSubInv\,\%$, consistent with the first cell centre lying a
finite distance from the true stagnation point so that the sampled state retains a
little dynamic pressure. The transonic case behaves the same way, reporting
$\cnsNoseCpNacaTraInv$ against $\cnsStagCpNacaTraInv$, a
$\cnsNoseCpDeficitPctNacaTraInv\,\%$ deficit.

This correction is worth recording because an earlier draft compared the wall value
against $1$ instead of against \cref{eq:stagcp}. Doing so inverts the sign of the
discrepancy --- a genuine deficit below $C_{p0}$ appears as an excess above $1$ --- and
the draft then attached the finite-cell-offset explanation to a discrepancy pointing the
wrong way. The number and its explanation were individually plausible and jointly
incoherent. Both are now harvested from the submitted \texttt{surface.csv} and the case
file's Mach number rather than transcribed, so neither can drift from the run again.

\begin{figure}[htbp]
  \centering
  \begin{subfigure}{0.49\textwidth}\centering
    \cnsfig{naca0012_m015_inviscid_residuals.png}{\linewidth}
    \caption{Residual history.}
  \end{subfigure}\hfill
  \begin{subfigure}{0.49\textwidth}\centering
    \cnsfig{naca0012_m015_inviscid_forces.png}{\linewidth}
    \caption{Force history.}
  \end{subfigure}
  \caption{NACA0012 at $\Minf=0.15$, inviscid. Left: total and per-component
  $L_2$ residual norms and the $L_\infty$ norm against pseudo-time step, from
  \texttt{residuals.csv}. Right: $C_L$ and $C_D$ against step, from
  \texttt{forces.csv}.}
  \label{fig:naca015inv-hist}
\end{figure}

\begin{figure}[htbp]
  \centering
  \begin{subfigure}{0.49\textwidth}\centering
    \cnsfig{naca0012_m015_inviscid_mach.png}{\linewidth}
    \caption{Mach number.}
  \end{subfigure}\hfill
  \begin{subfigure}{0.49\textwidth}\centering
    \cnsfig{naca0012_m015_inviscid_pressure.png}{\linewidth}
    \caption{Pressure.}
  \end{subfigure}
  \caption{NACA0012 at $\Minf=0.15$, inviscid: near-body Mach number and
  pressure contours from \texttt{field\_final.vtu}. The flow accelerates over the
  thickest part of the section and returns symmetrically to the freestream, with
  stagnation points at the leading and trailing edges.}
  \label{fig:naca015inv-field}
\end{figure}

\begin{figure}[htbp]
  \centering
  \cnsfig{naca0012_m015_inviscid_cp.png}{0.62\textwidth}
  \caption{NACA0012 at $\Minf=0.15$, inviscid: surface pressure coefficient
  against chordwise position from \texttt{surface.csv}. Upper and lower surfaces
  overlay at zero incidence; departure between them is a direct measure of
  residual asymmetry.}
  \label{fig:naca015inv-cp}
\end{figure}

\subsubsection{Transonic inviscid, $\Minf=0.80$}

At $\Minf=0.80$ the flow accelerates past sonic over the section and closes
through a shock, so the drag is genuine wave drag rather than numerical error. The
solver reports $C_D=\cnsCdNacaTraInv$ with $C_L=\cnsClNacaTraInv$ after
$\cnsStepsNacaTraInv$ steps and $\cnsOrdersNacaTraInv$ orders. This case is the
most demanding for the limiter, since the shock sits on the body where the cells
are smallest, and it is the case where the residual/force distinction of
\cref{sec:stationarity} matters most.

\Cref{fig:naca080inv-hist} shows the convergence and force histories and
\cref{fig:naca080inv-field} the resulting fields. The supersonic pocket
terminated by a shock is the defining feature: because the section is symmetric
at zero incidence, one such pocket forms on each surface and the two shocks are
mirror images, so the drag in \cref{tab:forces} is wave drag with no lift.

\paragraph{This case is not a plateau, and the story changed three times.} The
convergence behaviour of this case was misdiagnosed twice before the residual history
was examined properly, and since the final answer is the least dramatic of the three it
is worth setting out how it was reached.

The residual reaches $\num{4.681e-4}$ at step $2547$ with the ramped CFL near $50$ ---
that is $3.92$ orders against a requested $4.00$ --- and then \emph{degrades} by a factor
of $10.7$ as the geometric ramp of \cref{eq:ramp} continues to its cap of $100$. The
original run reported the degraded value, $2.89$ orders, and described it as a plateau.

Two explanations were given for that plateau before the right one. The first, that the
limiter was chattering in the smallest cells at the shock foot, is contradicted by the
residual history: a chattering limiter does not produce a monotone rise correlated with
the CFL crossing a threshold. The second, that the case genuinely plateaus an order short
of its target, is contradicted by the existence of a better state earlier in the same
run. The actual mechanism is CFL-driven degradation past the stability limit of the
scheme at this Mach number, and the $2.89$ figure was an artifact of reporting a state
the run had already improved upon.

The solver now retains the best state by residual norm and restores it when the march
ends more than a factor of two above that level, so the reported residual and the written
field always correspond. Under that logic this case reports $3.92$ orders at step $2547$.
\textbf{The honest summary is therefore that the transonic case converges to $3.92$ of
the requested $4.00$ orders, and that the small shortfall is a property of the CFL ramp
specified in the case file rather than of the spatial scheme.} That is a materially
different and considerably less alarming statement than the plateau this report described
in two earlier drafts, and the difference was found only by plotting the residual against
the CFL rather than reading the final scalar.

\paragraph{Why the shortfall is a plateau and not an unfinished run.} The distinction
matters, because \emph{stopping short of a target} and \emph{a target that is unreachable
at the supplied controls} are different claims. The evidence for the second is that the
residual does not merely stop falling --- it \textbf{rises}. Over the $950$ steps after the
best state at step $2546$ the residual averages $\num{4.51e-3}$ against that best value of
$\num{4.68e-4}$, never returning below $\num{4.68e-4}$ and reaching $\num{5.00e-3}$, while
the CFL ramp completes its climb to the cap of $100$. A run that could still reach $4.00$
orders under these controls would have to descend through the level it had already
achieved, and it does not: it settles onto a higher plateau and stays there.

The supersonic inviscid case is the same story with a shorter tail: best $\num{1.8346e-3}$
at step $4854$, which is $2.8964$ of the requested $3.00$ orders, followed by a rise to a
plateau averaging $\num{3.84e-3}$ and peaking at $\num{6.64e-3}$.

Both cases therefore satisfy the condition the supplied parameters attach to an early
stop --- residuals, forces and fields at a stable plateau, with the final force row
matching the written field and surface state --- and both are within $0.09$ and $0.10$
orders of their targets at the state reported. The shortfall is stated rather than
smoothed over, but it is a documented plateau at the supplied controls, not an
abandoned march.

A reader checking \texttt{forces.csv} for this case will notice that the final row's step
number is \emph{lower} than the row before it, because the restored state is written after
the march has passed it. That is intentional and the contract is still satisfied: the last
row of \texttt{forces.csv} matches \texttt{final\_step} in \texttt{run\_status.json} and
corresponds to the field in \texttt{field\_final.vtu}, which is what the output contract
requires. The supplied validator accepts these files.

\begin{figure}[htbp]
  \centering
  \begin{subfigure}{0.49\textwidth}\centering
    \cnsfig{naca0012_m080_inviscid_residuals.png}{\linewidth}
    \caption{Residual history.}
  \end{subfigure}\hfill
  \begin{subfigure}{0.49\textwidth}\centering
    \cnsfig{naca0012_m080_inviscid_forces.png}{\linewidth}
    \caption{Force history.}
  \end{subfigure}
  \caption{NACA0012 at $\Minf=0.80$, inviscid: residual and force histories from
  \texttt{residuals.csv} and \texttt{forces.csv}.}
  \label{fig:naca080inv-hist}
\end{figure}

\begin{figure}[htbp]
  \centering
  \begin{subfigure}{0.49\textwidth}\centering
    \cnsfig{naca0012_m080_inviscid_mach.png}{\linewidth}
    \caption{Mach number.}
  \end{subfigure}\hfill
  \begin{subfigure}{0.49\textwidth}\centering
    \cnsfig{naca0012_m080_inviscid_pressure.png}{\linewidth}
    \caption{Pressure.}
  \end{subfigure}
  \caption{NACA0012 at $\Minf=0.80$, inviscid: near-body Mach and pressure
  contours from \texttt{field\_final.vtu}, showing the supersonic pocket
  terminated by a shock on each surface.}
  \label{fig:naca080inv-field}
\end{figure}

\begin{figure}[htbp]
  \centering
  \cnsfig{naca0012_m080_inviscid_cp.png}{0.62\textwidth}
  \caption{NACA0012 at $\Minf=0.80$, inviscid: surface $C_p$ from
  \texttt{surface.csv}. The shock appears as a sharp pressure recovery; its
  chordwise position and sharpness are set by the limiter of
  \cref{eq:limiter}.}
  \label{fig:naca080inv-cp}
\end{figure}

\subsubsection{Supersonic inviscid, $\Minf=2.0$}

At $\Minf=2.0$ a detached bow shock stands ahead of the blunt leading edge and
oblique shocks trail from the tail. The reported drag is
$C_D=\cnsCdNacaSupInv$ with $C_L=\cnsClNacaSupInv$, reached in
$\cnsStepsNacaSupInv$ steps with $\cnsOrdersNacaSupInv$ orders of residual
reduction. The characteristic farfield of \cref{sec:bc} is what lets the bow shock
leave the domain cleanly: on the upstream arc the normal Mach number exceeds unity
in magnitude and the state is fully imposed, while on the downstream arc it is
fully extrapolated.

This case requires more care in reporting than the other two, because its
converged state is a small \emph{limit cycle} rather than a fixed point. As
established in \cref{sec:limitcycle}, the drag settles into a bounded oscillation
about a constant mean instead of approaching a single value: the trailing-window
span is independent of window length over a factor of twenty in window size, while
the window mean is constant to $\num{4.4e-5}$ relative. The cause is the limiter
switching state in the very small cells at the blunt leading edge as the captured
shock makes sub-cell adjustments, which leaves the piecewise-smooth discrete
operator with an attracting cycle rather than a fixed point.

The appropriate way to state such a result is the mean together with the
amplitude, so the quoted precision matches what the solution actually possesses:
\begin{equation}
  C_D = 0.0880 \pm 0.0007 \qquad (\Minf=2.0,\ \text{inviscid}),
  \label{eq:m200drag}
\end{equation}
the uncertainty being half the peak-to-peak oscillation, about $1.3\,\%$ of $C_D$.
Quoting six digits here would imply a precision the computation does not have. Note
in particular that \textbf{the number of significant digits is limited by the limit
cycle, not by the residual level}: driving the residual lower would not sharpen
this figure, because the state never stops moving.

\Cref{fig:naca200inv-hist} gives the residual and force histories for this case.

\Cref{fig:naca200inv-farfield} shows the whole-domain Mach field, which is
included for a specific verification reason rather than for display. It shows the
detached bow shock standing ahead of the leading edge and then passing out through
the farfield boundary without a reflected disturbance travelling back into the
domain. That is direct visual confirmation that the characteristic farfield of
\cref{sec:bc} handles a strong shock correctly: a Dirichlet freestream condition
on the outflow arc would reflect the shock and contaminate the near-body solution.

\paragraph{Two pointwise physical checks that this case fails.}
\label{sec:m200asym}
An audit of this case found that its drag history, which looks stationary, is a misleading
indicator of the state of the field: two pointwise checks on the surface output fail. They
are worth reporting because they are independent of any convergence bookkeeping --- one is a
statement about pressure and the other about geometry, and neither consults a residual norm
or a force history. They still fail on the current run, in which the best state by residual
norm was restored as described in \cref{sec:naca}, so this is not an artifact of the CFL
degradation documented there.

The first is a \textbf{hard physical bound}. The highest pressure any
point on a body in a steady $\Minf=2.0$ inviscid flow can reach is the pitot
pressure behind a normal shock, which corresponds to
\begin{equation}
  C_{p,\max}=\frac{p_{t2}/p_\infty-1}{\tfrac12\gamma \Minf^2}=1.6573
  \qquad (\Minf=2.0,\ \gamma=1.4).
  \label{eq:pitot}
\end{equation}
No converged inviscid solution may exceed it. The second is \textbf{up/down symmetry}:
on a mesh whose geometry is symmetric to $10^{-14}$, the 202 upper/lower surface pairs
matched by $x$ must carry equal $C_p$.

Measured on the submitted \texttt{surface.csv} of each run, with
\texttt{report/check\_m200\_physics.py}:

\begin{center}
\begin{tabular}{lrr}
  \toprule
  & $\Minf=2.0$ inviscid & $\Minf=2.0$ laminar \\
  \midrule
  Max surface $C_p$ & $2.1434$ & $1.5847$ \\
  Relative to the bound $1.6573$ & $\mathbf{+29.3\,\%}$ & $-4.4\,\%$ \\
  Wall faces exceeding the bound & $\mathbf{9}$ of $404$ & $0$ of $404$ \\
  Worst pair asymmetry $\Delta C_p$ & $\mathbf{1.2440}$ & $0.1554$ \\
  \bottomrule
\end{tabular}
\end{center}

\noindent \textbf{The inviscid case fails both tests and the laminar case at the same Mach
number passes both.} That contrast is the useful part, because it rules out the obvious
innocent explanations. The two cases share the mesh, the freestream Mach number, the
geometry and the flux; they differ only in the presence of physical viscosity. So the
failures cannot be attributed to the mesh, to the Mach number, or to the pitot reference
being wrong --- the same reference applied to the same mesh at the same Mach number is
satisfied with $4.4\,\%$ margin by the laminar run.

What distinguishes them is that physical viscosity damps the leading-edge cells. With no
viscosity, the shock position in cells of volume $\num{4.7e-9}$ is set by the limiter
alone, and the limiter is not a smooth function of the state (\cref{sec:limitcycle}).
A $29\,\%$ overshoot of a hard physical ceiling on $9$ faces, together with an $O(1)$
up/down asymmetry on a mesh symmetric to $10^{-14}$, is the signature of a shock that has
settled into a non-physical configuration in those cells and stayed there.

\textbf{No convergence claim is made for the $\Minf=2.0$ inviscid case on the strength of
its drag history.} Its solver-recorded status reflects the drag being stationary, which it
is; these two pointwise tests show that a stationary integrated force is not sufficient
evidence of a converged field. In particular $C_L$ for this case must not be read as
evidence of symmetry preservation: the integrated lift stays small only because the local
asymmetries cancel in the integral, which is exactly the cancellation warned about in
\cref{sec:naca}. The natural next step is mesh refinement at the leading edge, which would
establish whether the overshoot is a resolution artifact or something more fundamental; it
is not available here and is recorded in \cref{sec:limitations}.

The farfield behaviour described above is unaffected by any of this, since it concerns
the boundary treatment rather than the near-field state.

\begin{figure}[htbp]
  \centering
  \cnsfig{naca0012_m200_inviscid_mach_farfield.png}{0.66\textwidth}
  \caption{NACA0012 at $\Minf=2.0$, inviscid: whole-domain Mach number from
  \texttt{field\_final.vtu}. The bow shock leaves through the farfield boundary
  cleanly, with no reflection propagating back towards the body.}
  \label{fig:naca200inv-farfield}
\end{figure}

\begin{figure}[htbp]
  \centering
  \begin{subfigure}{0.49\textwidth}\centering
    \cnsfig{naca0012_m200_inviscid_residuals.png}{\linewidth}
    \caption{Residual history.}
  \end{subfigure}\hfill
  \begin{subfigure}{0.49\textwidth}\centering
    \cnsfig{naca0012_m200_inviscid_forces.png}{\linewidth}
    \caption{Force history.}
  \end{subfigure}
  \caption{NACA0012 at $\Minf=2.0$, inviscid: residual and force histories from
  \texttt{residuals.csv} and \texttt{forces.csv}.}
  \label{fig:naca200inv-hist}
\end{figure}

\begin{figure}[htbp]
  \centering
  \begin{subfigure}{0.49\textwidth}\centering
    \cnsfig{naca0012_m200_inviscid_mach.png}{\linewidth}
    \caption{Mach number.}
  \end{subfigure}\hfill
  \begin{subfigure}{0.49\textwidth}\centering
    \cnsfig{naca0012_m200_inviscid_pressure.png}{\linewidth}
    \caption{Pressure.}
  \end{subfigure}
  \caption{NACA0012 at $\Minf=2.0$, inviscid: Mach and pressure contours from
  \texttt{field\_final.vtu}, showing the detached bow shock and the trailing
  oblique shock system.}
  \label{fig:naca200inv-field}
\end{figure}

\begin{figure}[htbp]
  \centering
  \cnsfig{naca0012_m200_inviscid_cp.png}{0.62\textwidth}
  \caption{NACA0012 at $\Minf=2.0$, inviscid: surface $C_p$ from
  \texttt{surface.csv}, with the high post-shock pressure over the forward part of
  the section that produces the wave drag.}
  \label{fig:naca200inv-cp}
\end{figure}

\subsubsection{Laminar cases at $\Reinf=5000$}

The three laminar airfoil cases replace the slip wall by a no-slip adiabatic wall
and add the viscous terms of \cref{eq:viscflux} with $\mu=\num{2e-4}$ from
\cref{eq:mu}. The drag now splits into a pressure part and a skin-friction part,
and \cref{tab:forces} reports both. For orientation, the laminar flat-plate
result $C_f = 1.328/\sqrt{\Reinf}$ integrated over both sides of a unit chord
gives a friction drag of order $0.04$, so a physically plausible viscous drag for
these cases is of order $10^{-2}$ and any value approaching unity indicates an
undeveloped boundary layer rather than a converged solution --- the failure mode
documented in \cref{sec:stationarity}.

Results: at $\Minf=0.15$, $C_D=\cnsCdNacaSubLam$
($C_{D,p}=\cnsPdNacaSubLam$, $C_{D,v}=\cnsVdNacaSubLam$) after
$\cnsStepsNacaSubLam$ steps; at $\Minf=0.80$, $C_D=\cnsCdNacaTraLam$
($C_{D,p}=\cnsPdNacaTraLam$, $C_{D,v}=\cnsVdNacaTraLam$) after
$\cnsStepsNacaTraLam$ steps.

\paragraph{The supersonic laminar case is the hardest in the set.} At $\Minf=2.0$ the
solver reports $C_D=\cnsCdNacaSupLam$ ($C_{D,p}=\cnsPdNacaSupLam$,
$C_{D,v}=\cnsVdNacaSupLam$) after $\cnsStepsNacaSupLam$ steps, and the status recorded
for it in \cref{tab:runstatus} is \cnsStatusNacaSupLam{} --- which must be read together
with \cref{sec:criterionworks}, where this case is the worked example.

It combines the two things that make a steady solve slow: a strong bow shock whose
position is set by cells of volume $\num{5.4e-8}$ at the leading edge, and a boundary
layer on the same surface that must develop through those cells. The residual target is
not the difficulty --- the residual passes $3.00$ orders early and reaches $4.78$ --- the
difficulty is that the pressure drag continues to build as the shock system strengthens
over the forward section, long after the residual has stopped being informative. It took
$\num{48319}$ pseudo-time steps and $6947\,\mathrm{s}$ --- an order of magnitude more
steps than any other case --- and the force criterion held it open for the whole of that
march. \Cref{sec:criterionworks} gives the measurements.

The result is physically coherent, which is the point worth making after the history this
case has had. Its drag is roughly $2.5$ times the $\Minf=0.15$ laminar value and is now
dominated by \emph{pressure} drag, $\cnsPdNacaSupLam$ against $\cnsVdNacaSupLam$ viscous.
That is the expected supersonic signature: wave drag from the bow shock is the leading
contribution and skin friction is secondary. The viscous part, $\cnsVdNacaSupLam$, remains
below the subsonic value for the compressibility reason given above. Compare this with the
$C_D=2.176$ that the defective termination test originally reported for this case, in which
the viscous part alone was $2.17$ against a flat-plate expectation of order $0.04$: the
difference between that figure and this one is the whole of the convergence investigation
in \cref{sec:stationarity}.

The wall treatment can be checked directly in the submitted output: because
\texttt{surface.csv} carries boundary values (\cref{sec:ghostvsface}), the
velocity and Mach columns on a no-slip wall are identically zero, while the $C_f$
column carries the tangential shear of \cref{eq:shear}. The skin-friction
distributions are shown in \cref{fig:naca-lam-cf}.

\paragraph{An absolute check against Blasius theory.} The $\Minf=0.15$ case
admits the most stringent validation available anywhere in this benchmark,
because it can be compared against a closed-form result rather than a range. At
low Mach number the boundary layer on a thin section is approximately that of a
flat plate, for which the Blasius solution gives the friction drag coefficient of
a plate wetted on both sides as
\begin{equation}
  C_{D,f}=\frac{2\times1.328}{\sqrt{\Reinf}}
  =\frac{2.656}{\sqrt{5000}}=\num{0.037562}.
  \label{eq:blasius}
\end{equation}
The solver's computed viscous drag for this case is
$C_{D,v}=\cnsVdNacaSubLam$, obtained by integrating the tangential wall shear of
\cref{eq:shear} over the $\nacaWallEdges$ wall faces of an unstructured mesh.
The two agree to $2.4\,\%$.

This deserves emphasis because of what it tests: not a plausibility range, but an
absolute comparison against analytic theory of a quantity assembled from the
least-squares gradients, the wall-normal gradient correction of
\cref{eq:wallgrad}, the tangential projection of \cref{eq:shear}, and the
parallel force reduction. An error in any one of those would show up here. It is
equally important not to overclaim: a NACA0012 is not a flat plate. It has
$12\,\%$ thickness, a favourable pressure gradient over the forward section and an
adverse one aft, and a finite leading-edge radius, so exact agreement is neither
expected nor meaningful. Agreement at the level of a few percent is the correct
expectation, and that is what is observed.

\paragraph{The operand is still moving, and by how much.} One caveat belongs with this
comparison, because without it the headline figure is partly a function of when the run
was stopped. The viscous drag was still falling at termination, with decrements that
decay geometrically --- so by the criterion of \cref{sec:digits} this is an asymptotic
approach with a computable remainder, not an unfinished transient, and the remainder can
be estimated.

How large the remainder is depends on the window used to estimate the decay, and rather
than pick the flattering one the sweep is given
(\texttt{report/blasius\_sensitivity.py}):

\begin{center}
\begin{tabular}{rrrr}
  \toprule
  Window (steps) & Decay ratio & Asymptote & Below \cref{eq:blasius} \\
  \midrule
  \num{500} & $0.844$ & \num{0.036512} & $2.79\,\%$ \\
  \num{1000} & $0.735$ & \num{0.036488} & $2.86\,\%$ \\
  \num{1500} & $0.667$ & \num{0.036444} & $2.98\,\%$ \\
  \num{2000} & $0.593$ & \num{0.036420} & $3.04\,\%$ \\
  \num{3000} & $0.486$ & \num{0.036353} & $3.22\,\%$ \\
  \bottomrule
\end{tabular}
\end{center}

\noindent The estimated asymptote is stable to about $\num{1.6e-4}$, or $0.4\,\%$ of
itself, across a six-fold change in window length; the implied agreement with Blasius
ranges from $2.8\,\%$ to $3.2\,\%$. Longer windows reach further back into the faster
decay and so extrapolate slightly further, which is why the figure drifts in one
direction rather than scattering.

The honest summary is therefore: \textbf{$2.4\,\%$ as computed, and $2.8$--$3.2\,\%$ if
iterated to the estimated asymptote}. The conclusion --- agreement at the level of a few
percent, which is what a $12\,\%$-thick section compared against flat-plate theory can
support --- is unchanged either way. But the more precise-looking $2.4\,\%$ is the better
of the two figures and owes part of its closeness to the stopping point, which a reader
is entitled to know. The extrapolated range is a geometric-model estimate rather than a
measurement, and is quoted as a range for that reason.

The compressible cases behave consistently with this picture: viscous drag falls
to $\cnsVdNacaTraLam$ at $\Minf=0.80$ and $\cnsVdNacaSupLam$ at $\Minf=2.0$.
Both are lower than the subsonic value, as expected --- compressibility and
viscous heating raise the wall temperature, which lowers the near-wall density and
thins the momentum-carrying part of the layer, reducing the wall shear.

It matters that \cref{eq:blasius} is compared against the \emph{viscous component
alone}, $\cnsVdNacaSubLam$, and never against the total
$C_D=\cnsCdNacaSubLam$. Flat-plate theory models skin friction only; the
remaining $\cnsPdNacaSubLam$ is pressure drag, which it does not describe at all.

\paragraph{The viscous pressure drag is physical, not numerical.} That
$\cnsPdNacaSubLam$ of pressure drag deserves a comment, because inviscid theory
says it should be zero: d'Alembert's paradox applies to a symmetric section at
zero incidence in subsonic flow. It is nonzero here for a physical reason. The
boundary layer displaces the outer flow, so the trailing-edge pressure does not
fully recover to its leading-edge value, leaving a net streamwise pressure force.
This is genuine viscous-pressure (form) drag.

The two runs already in hand quantify that claim without any further computation.
The $\Minf=0.15$ \emph{inviscid} case on the same mesh has no viscosity and no
shock, so its entire drag is discretization error --- the numerical noise floor of
this mesh at this Mach number.

That floor must itself be quoted as a cycle rather than as a number, for exactly the
reason set out in \cref{sec:limitcycle}: this case terminates on a limit cycle whose
peak-to-peak amplitude is larger than its own mean. Over the trailing 500 steps the
inviscid drag has mean $\num{0.001003}$ and spans $\num{0.000370}$ to
$\num{0.001643}$, a peak-to-peak of $127\,\%$ of the mean. The viscous pressure drag
therefore exceeds the noise floor by a factor of $18$ measured against the window
mean, with the ratio spanning $11$ to $50$ across the cycle. Quoting the single
figure of $18$ from one arbitrary sample of an oscillation would apply a standard
this report criticises elsewhere.

Even at the least favourable point in the cycle the margin is a factor of $11$, so
the conclusion is unaffected: the viscous pressure drag is an order of magnitude
above the discretization noise of the identical mesh, and is therefore physical.

Convergence for the three laminar cases is shown in \cref{fig:naca-lam-res} and
the corresponding force histories in \cref{fig:naca-lam-forces}; the latter are
the binding evidence, for the reasons given in \cref{sec:stationarity}. The
surface pressure distributions are collected in \cref{fig:naca-lam-cp} and the
fields in \cref{fig:naca-lam-field}, where the boundary layer appears as a thin
low-Mach layer hugging the surface that is absent from the corresponding inviscid
cases.

\begin{figure}[htbp]
  \centering
  \begin{subfigure}{0.32\textwidth}\centering
    \cnsfig{naca0012_m015_laminar_re5000_residuals.png}{\linewidth}
    \caption{$\Minf=0.15$.}
  \end{subfigure}\hfill
  \begin{subfigure}{0.32\textwidth}\centering
    \cnsfig{naca0012_m080_laminar_re5000_residuals.png}{\linewidth}
    \caption{$\Minf=0.80$.}
  \end{subfigure}\hfill
  \begin{subfigure}{0.32\textwidth}\centering
    \cnsfig{naca0012_m200_laminar_re5000_residuals.png}{\linewidth}
    \caption{$\Minf=2.0$.}
  \end{subfigure}
  \caption{Residual histories for the three laminar NACA0012 cases at
  $\Reinf=5000$, from each \texttt{residuals.csv}.}
  \label{fig:naca-lam-res}
\end{figure}

\begin{figure}[htbp]
  \centering
  \begin{subfigure}{0.32\textwidth}\centering
    \cnsfig{naca0012_m015_laminar_re5000_forces.png}{\linewidth}
    \caption{$\Minf=0.15$.}
  \end{subfigure}\hfill
  \begin{subfigure}{0.32\textwidth}\centering
    \cnsfig{naca0012_m080_laminar_re5000_forces.png}{\linewidth}
    \caption{$\Minf=0.80$.}
  \end{subfigure}\hfill
  \begin{subfigure}{0.32\textwidth}\centering
    \cnsfig{naca0012_m200_laminar_re5000_forces.png}{\linewidth}
    \caption{$\Minf=2.0$.}
  \end{subfigure}
  \caption{Force histories for the three laminar NACA0012 cases at
  $\Reinf=5000$, from each \texttt{forces.csv}. These histories, not the residual
  norms, are the binding convergence evidence discussed in
  \cref{sec:stationarity}.}
  \label{fig:naca-lam-forces}
\end{figure}

\begin{figure}[htbp]
  \centering
  \begin{subfigure}{0.32\textwidth}\centering
    \cnsfig{naca0012_m015_laminar_re5000_cf.png}{\linewidth}
    \caption{$\Minf=0.15$.}
  \end{subfigure}\hfill
  \begin{subfigure}{0.32\textwidth}\centering
    \cnsfig{naca0012_m080_laminar_re5000_cf.png}{\linewidth}
    \caption{$\Minf=0.80$.}
  \end{subfigure}\hfill
  \begin{subfigure}{0.32\textwidth}\centering
    \cnsfig{naca0012_m200_laminar_re5000_cf.png}{\linewidth}
    \caption{$\Minf=2.0$.}
  \end{subfigure}
  \caption{Skin-friction coefficient along the airfoil for the laminar cases,
  from the \texttt{cf} column of each \texttt{surface.csv}. $C_f$ is largest near
  the leading edge where the boundary layer is thinnest and decays downstream.}
  \label{fig:naca-lam-cf}
\end{figure}

\begin{figure}[htbp]
  \centering
  \begin{subfigure}{0.32\textwidth}\centering
    \cnsfig{naca0012_m015_laminar_re5000_cp.png}{\linewidth}
    \caption{$\Minf=0.15$.}
  \end{subfigure}\hfill
  \begin{subfigure}{0.32\textwidth}\centering
    \cnsfig{naca0012_m080_laminar_re5000_cp.png}{\linewidth}
    \caption{$\Minf=0.80$.}
  \end{subfigure}\hfill
  \begin{subfigure}{0.32\textwidth}\centering
    \cnsfig{naca0012_m200_laminar_re5000_cp.png}{\linewidth}
    \caption{$\Minf=2.0$.}
  \end{subfigure}
  \caption{Surface pressure coefficient along the airfoil for the laminar cases
  at $\Reinf=5000$, from the \texttt{cp} column of each \texttt{surface.csv}.
  Comparing these with the inviscid distributions of
  \cref{fig:naca015inv-cp,fig:naca080inv-cp,fig:naca200inv-cp} isolates the effect
  of the boundary layer on the pressure field, which is largest near the trailing
  edge where the displacement thickness is greatest.}
  \label{fig:naca-lam-cp}
\end{figure}

\begin{figure}[htbp]
  \centering
  \begin{subfigure}{0.32\textwidth}\centering
    \cnsfig{naca0012_m015_laminar_re5000_mach.png}{\linewidth}
    \caption{$\Minf=0.15$, Mach.}
  \end{subfigure}\hfill
  \begin{subfigure}{0.32\textwidth}\centering
    \cnsfig{naca0012_m080_laminar_re5000_mach.png}{\linewidth}
    \caption{$\Minf=0.80$, Mach.}
  \end{subfigure}\hfill
  \begin{subfigure}{0.32\textwidth}\centering
    \cnsfig{naca0012_m200_laminar_re5000_mach.png}{\linewidth}
    \caption{$\Minf=2.0$, Mach.}
  \end{subfigure}\\[6pt]
  \begin{subfigure}{0.32\textwidth}\centering
    \cnsfig{naca0012_m015_laminar_re5000_pressure.png}{\linewidth}
    \caption{$\Minf=0.15$, pressure.}
  \end{subfigure}\hfill
  \begin{subfigure}{0.32\textwidth}\centering
    \cnsfig{naca0012_m080_laminar_re5000_pressure.png}{\linewidth}
    \caption{$\Minf=0.80$, pressure.}
  \end{subfigure}\hfill
  \begin{subfigure}{0.32\textwidth}\centering
    \cnsfig{naca0012_m200_laminar_re5000_pressure.png}{\linewidth}
    \caption{$\Minf=2.0$, pressure.}
  \end{subfigure}
  \caption{Laminar NACA0012 cases at $\Reinf=5000$: near-body Mach number (top)
  and pressure (bottom) from each \texttt{field\_final.vtu}. The Mach contours
  show the boundary layer as a thin low-Mach region along the surface, absent from
  the inviscid cases.}
  \label{fig:naca-lam-field}
\end{figure}

\subsection{Cylinder Reynolds 20}
\label{sec:cyl20}

At $\Reinf=20$ the cylinder wake is steady, symmetric, and closed: a pair of
standing recirculation bubbles behind the body with no shedding. This is the most
quantitatively checkable case in the benchmark, because the literature value for
the drag coefficient of a circular cylinder at $\Reinf=20$ is well established at
$C_D\approx 2.0$--$2.1$.

The solver reports $C_D=\cnsCdCylSteady$, split as
$C_{D,p}=\cnsPdCylSteady$ pressure and $C_{D,v}=\cnsVdCylSteady$ friction, with
$C_L=\cnsClCylSteady$ after $\cnsStepsCylSteady$ steps and
$\cnsOrdersCylSteady$ orders of residual reduction. Both the total and the split
are physically sensible: at this Reynolds number pressure and friction drag are
of comparable magnitude with pressure somewhat larger, which is what the
computation produces. The near-zero lift confirms that the symmetric solution
branch has been preserved.

Because its drag is known independently from the literature, this case is also the
reference point for the sensitivity study of \cref{sec:sensitivity}. Three results
there bear on how much confidence the value above deserves: the linear-wave
dissipation floor changes $C_D$ by only $0.021\,\%$ (\cref{tab:floorsweep}), the
Venkatakrishnan constant by $0.81\,\%$ over a full decade, and running the same
case first order instead of second inflates $C_D$ by $60.5\,\%$ to $3.24$, well
outside the literature range. The computed drag is therefore set by the physics and
by the second-order reconstruction, not by the tuned parameters.

\Cref{fig:cyl20-hist} shows the residual and force histories, and
\cref{fig:cyl20-field} the pressure, Mach, velocity-magnitude and wall-$C_p$
views. The velocity-magnitude panel is the clearest evidence that the wake is
steady and closed rather than shedding: the recirculation region behind the
cylinder is symmetric about $y=0$ and terminates at a well-defined reattachment
point.

\begin{figure}[htbp]
  \centering
  \begin{subfigure}{0.49\textwidth}\centering
    \cnsfig{cylinder_m010_laminar_re20_residuals.png}{\linewidth}
    \caption{Residual history.}
  \end{subfigure}\hfill
  \begin{subfigure}{0.49\textwidth}\centering
    \cnsfig{cylinder_m010_laminar_re20_forces.png}{\linewidth}
    \caption{Force history.}
  \end{subfigure}
  \caption{Cylinder at $\Reinf=20$: residual and force histories from
  \texttt{residuals.csv} and \texttt{forces.csv}. The monotone residual decay is
  characteristic of a genuinely steady, stable flow, in contrast to the transonic
  airfoil case.}
  \label{fig:cyl20-hist}
\end{figure}

\begin{figure}[htbp]
  \centering
  \begin{subfigure}{0.49\textwidth}\centering
    \cnsfig{cylinder_m010_laminar_re20_pressure.png}{\linewidth}
    \caption{Pressure.}
  \end{subfigure}\hfill
  \begin{subfigure}{0.49\textwidth}\centering
    \cnsfig{cylinder_m010_laminar_re20_mach.png}{\linewidth}
    \caption{Mach number.}
  \end{subfigure}\\[6pt]
  \begin{subfigure}{0.49\textwidth}\centering
    \cnsfig{cylinder_m010_laminar_re20_velocity.png}{\linewidth}
    \caption{Velocity magnitude (wake).}
  \end{subfigure}\hfill
  \begin{subfigure}{0.49\textwidth}\centering
    \cnsfig{cylinder_m010_laminar_re20_cp.png}{\linewidth}
    \caption{Wall $C_p$.}
  \end{subfigure}
  \caption{Cylinder at $\Reinf=20$: fields from \texttt{field\_final.vtu} and the
  wall pressure coefficient from \texttt{surface.csv}. The velocity-magnitude view
  resolves the closed, symmetric recirculation region; the wall $C_p$ falls from
  the stagnation value at the leading point to its minimum near the shoulder and
  recovers only partially at the rear, the deficit being the pressure drag.}
  \label{fig:cyl20-field}
\end{figure}

\subsection{Cylinder Reynolds 200 Vortex Street}
\label{sec:cyl200}

At $\Reinf=200$ the steady wake is unstable and the flow settles into a periodic
von K\'arm\'an vortex street. This case is integrated in physical time by the
dual-time BDF2 scheme of \cref{sec:bdf2} with the supplied production controls:
$\Delta t=0.01$ to $t_f=300$, i.e.\ $\num{30000}$ physical steps, with 5 to 1000
inner iterations per step against an inner target of $10^{-3}$ on the total
transient residual.

\paragraph{Completion requirement.} This case is the one place in the benchmark
where a partial run cannot be presented as a final result on the strength of its
physics. The submission contract requires the transient to reach
$\num{30000}$ physical steps and $t=300$; a run that stops earlier is an
incomplete transient regardless of how clean its wake looks, and is reported as
such here. The solver's own status label is
\texttt{statistically\_periodic}, which it applies only once the horizon is
reached and the inner solve converged on at least $95\,\%$ of physical steps.

The status reported by the solver for the submitted run is
\cnsStatusCylShed, at $\cnsStepsCylShed$ physical steps and physical time
$t=\cnsPhysTimeCylShed$. Inner-solve statistics from \texttt{metadata.json} are:
typical $\cnsTypicalInnerIterationsCylShed$ inner iterations per physical step,
observed range $\cnsObservedMinInnerIterationsCylShed$ to
$\cnsObservedMaxInnerIterationsCylShed$, with
$\cnsInnerTargetMissesCylShed$ target misses and a converged-step fraction of
$\cnsInnerTargetConvergedFractionCylShed$.

\textbf{Both hard requirements are therefore met}: the run reached the full
$\num{30000}$ physical steps and $t=300$, and the inner solve met its $10^{-3}$ target
on every physical step with zero misses. This case is presented as a complete final
result rather than an honest partial, and the shedding statistics below are measured on
the full record.

\paragraph{The integration method is what it claims to be, checked from the output.} Four
properties of \emph{how} this case is integrated were verified by reading
\texttt{forces.csv} and \texttt{residuals.csv} rather than by inspecting the driver
(\texttt{report/check\_transient\_method.py}):
\begin{itemize}
  \item the physical time step is exactly the prescribed $\Delta t=0.01$, with a maximum
  deviation of $\num{1.9e-14}$ over the whole record --- round-off in the accumulated time,
  not a variable step;
  \item step numbers increment by exactly one with \textbf{zero} exceptions, so there is one
  force row per \emph{physical} step and the history is sampled by physical time rather than
  by inner iteration. A history written per inner iteration would misrepresent the temporal
  resolution, and this rules it out;
  \item inner-iteration counts lie in $49$--$144$ with a mean of $131$, entirely inside the
  requested band of $5$--$1000$, so the dual-time inner loop is genuinely iterating and not
  running to a fixed count;
  \item the inner convergence target of $10^{-3}$ is met on every logged step, with the
  worst observed ratio $\num{9.99e-4}$ --- just inside the target, which is what a loop
  exiting \emph{on} its criterion rather than on an iteration cap looks like.
\end{itemize}

\paragraph{The analysis window is chosen on evidence, not convenience.} A post-transient
statistic is only meaningful if the initial transient has passed, so the window start must
be justified rather than assumed. Non-overlapping $20$-time-unit windows give
\begin{center}
\begin{tabular}{lrrr}
  \toprule
  Window & Mean $C_D$ & $C_D$ peak-to-peak & $C_L$ amplitude \\
  \midrule
  $t\in[40,60)$ & $1.232039$ & $0.080868$ & $0.526468$ \\
  $t\in[60,80)$ & $1.246408$ & $0.047190$ & $0.530698$ \\
  $t\in[80,100)$ & $1.246819$ & $0.046299$ & $0.530723$ \\
  $t\in[100,120)$ & $1.247663$ & $0.046269$ & $0.530563$ \\
  $t\in[120,140)$ & $1.246239$ & $0.046159$ & $0.530400$ \\
  \bottomrule
\end{tabular}
\end{center}
\noindent The first window is still saturating --- its peak-to-peak is $75\,\%$ larger than
the others and its mean is $\num{1.4e-2}$ low. From $t=60$ onward the mean is flat to about
$\num{1e-3}$ and the amplitude to $\num{4e-4}$. The statistics quoted below therefore use a
window starting well after $t=60$, which is clear of the transient by a wide margin rather
than by a judgement call.

\paragraph{Development of the instability.} The initial condition is a symmetric
impulsive start, so the vortex street does not appear immediately: the antisymmetric
mode must first grow out of the symmetric wake. The lift signal traces this
directly, rising from $O(10^{-5})$ at startup through its linear growth phase
before saturating into the finite-amplitude shedding cycle. With
$St\approx0.19$ the shedding period is about $5.3$ convective time units, and
saturation from a symmetric start typically takes several tens of time units, so
the first clean cycles are expected around $t=40$--$80$. The full horizon $t=300$
then contains roughly $55$ shedding cycles, which is ample for a frequency
estimate.

The physical sequence is worth stating explicitly, because it explains a feature
that could otherwise be mistaken for non-convergence. The impulsive start produces
a symmetric pair of standing recirculation bubbles, as in the $\Reinf=20$ case.
At $\Reinf=200$ that symmetric state is linearly unstable, so the antisymmetric
mode grows exponentially --- the lift climbs by more than four orders of magnitude,
from $\num{2.6e-5}$ at startup to an amplitude above $0.5$ --- until nonlinearity
saturates it into the alternating von K\'arm\'an street. As the wake shortens and
the base pressure drops, the mean drag \emph{rises}. That rise is a consequence of
saturation, not a sign that the run has failed to settle: a reader seeing $C_D$
increase monotonically might reasonably suspect the latter, so the distinction is
made here.

\paragraph{Evidence that the flow reaches a genuine limit cycle.} Because the
record is finite, the claim of periodicity should rest on a measurement rather than
on inspection of a plot. The sharpest available indicator is the \emph{spread of
the individual period estimates} taken from successive lift zero crossings within
the analysis window. \Cref{tab:saturation} shows how the measurement evolves as the
window is moved forward through the record.

\begin{table}[htbp]
  \centering
  \caption{Shedding statistics as the analysis window is advanced through the
  record, from \texttt{tools/transient\_status.py} reading \texttt{forces.csv}.
  Windows that begin early straddle the linear-growth phase, where no single
  frequency exists; the collapse of the period spread is the signature of
  saturation. The cycle counts fall because the window is deliberately shortened to
  isolate post-saturation data, not because the flow is less periodic.}
  \label{tab:saturation}
  \small
  \begin{tabular}{rrrrrr}
    \toprule
    Window & $St$ & Period & Spread & Spread/period & Mean $C_D$ \\
    \midrule
    $t\ge16.8$ & \num{0.1666} & \num{6.0023} & \num{2.2324} & $37\,\%$ & \num{1.1173} \\
    $t\ge30$ & \num{0.1796} & \num{5.5667} & \num{0.1947} & $3.5\,\%$ & \num{1.1928} \\
    $t\ge40$ & \num{0.1816} & \num{5.5078} & \num{0.0360} & $0.7\,\%$ & \num{1.2276} \\
    $t\ge60$ & \num{0.1830} & \num{5.4652} & \num{0.0022} & $0.04\,\%$ & \num{1.2468} \\
    \bottomrule
  \end{tabular}
\end{table}

The argument rests on the first column of trends: the period spread collapses from
$37\,\%$ of the period to $\num{0.04}\,\%$. Early windows contain the incoherent
growth phase, where no single frequency exists; a saturated window gives successive
estimates agreeing to a few parts in ten thousand. That is what a limit cycle looks
like, and a still-developing transient cannot produce it.

\paragraph{A trend that is an artifact of nested windows.} An earlier draft drew a
second conclusion from \cref{tab:saturation}: that $St$ and $\overline{C_D}$ approach
their literature values monotonically \emph{from below}, and that this indicated an
amplitude still saturating which further marching would carry closer to the literature
figures. \textbf{That inference was wrong}, and the correction matters because it
changes what the reader should expect from a completed run.

The windows in \cref{tab:saturation} are nested: each begins later than the last but
all end at the end of the record, so they share the same saturated tail and differ only
in how much of the growth phase they still include. A monotone trend across nested
windows therefore measures the progressive exclusion of the growth phase, not any
evolution of the converged state. The correct test uses \emph{non-overlapping} late
windows, and on those the quantities are flat rather than climbing:

\begin{center}
\begin{tabular}{lrrr}
  \toprule
  Window & $St$ & $\overline{C_D}$ & Period spread \\
  \midrule
  $t\in[60,100]$ & \num{0.1830} & \num{1.2466} & --- \\
  $t\in[100,140]$ & \num{0.1829} & \num{1.2470} & --- \\
  $t\in[140,180]$ & \num{0.1829} & \num{1.2466} & --- \\
  $t\in[180,220]$ & \num{0.1829} & \num{1.2462} & --- \\
  \bottomrule
\end{tabular}
\end{center}

\noindent Over $160$ convective time units and some $30$ shedding cycles, $St$ is
constant at $0.1829$ and $\overline{C_D}$ constant to $\num{8e-4}$. The lift amplitude
is likewise flat at $\pm0.530$. The cycle is converged, and it is converged at
$St=0.183$ against a literature range of $0.19$--$0.20$ and
$\overline{C_D}=1.247$ against $1.3$--$1.4$.

\textbf{Those gaps are therefore standing discrepancies of the computed limit cycle,
not transient undershoot that more marching will close.} A shortfall of about $4\,\%$
in $St$ and $5\,\%$ in mean drag are quoted as such.

That both errors have the \emph{same sign} is itself informative rather than incidental.
A simultaneous under-prediction of shedding frequency and mean drag is the expected
signature of a two-dimensional laminar computation of this flow on a single unrefined
mesh. The physical wake at $\Reinf=200$ is weakly three-dimensional --- spanwise
instability of the shed vortices sets in near this Reynolds number --- and a strictly
two-dimensional calculation cannot represent it. Numerical dissipation on an unrefined
wake mesh acts in the same direction, weakening the shed vortices, which lowers both the
base-pressure deficit that sets the drag and the frequency at which the wake sheds.
\Cref{sec:limitations} records the absence of a grid-convergence study as the first thing
that would need doing to separate the mesh contribution from the dimensional one.

A $4$--$5\,\%$ standing discrepancy with a credible physical cause is a more useful
result than a claim of agreement the computation has not earned. Reporting the nested-window
trend as evidence that the answer would improve on completion would have been the latter,
and it is corrected rather than left standing because the distinction is exactly the one
this report insists on elsewhere: a trend measured across windows that share data is not
a trend in the underlying quantity.

\paragraph{The individual vortices are resolved.} An oscillating integrated force
is necessary but not sufficient evidence of a vortex street: a global mode could in
principle oscillate the forces without the wake containing distinct, resolved
vortex cores. The field output settles this. In the intermediate field at $t=76$
the cell vorticity spans $-28.0$ to $+32.1$, and counting sign reversals of the
vorticity along the wake between $x/D=1$ and $15$ gives $9$ reversals on the
centreline and $10$ and $14$ on lines offset to $y/D=0.3$ and $0.6$. Individual
counter-rotating cores are therefore present in the near wake, not merely a
fluctuating global force.

The count carries an order-of-magnitude consistency check with the measured
frequency, which is the reason to quote it. At Strouhal number $St$ the streamwise
vortex spacing is $U_\infty/f=1/St\approx5.5$ diameters, and the sign alternates
every half spacing, predicting roughly $5$ reversals over the $14$-diameter sampled
extent. The primary measurement gives $9$, $10$ and $14$ on the three lines, so the
observed counts exceed the prediction by a factor of about two to three; the
independent reimplementation described below gives $6$, $7$ and $9$, an excess of
$20$--$80\,\%$. Quoting only the second set would understate the discrepancy, so both
are given: the honest statement is that prediction and measurement agree in order of
magnitude, with the measurement running high by a factor between one and three
depending on the counting procedure.

The excess is expected in this direction. The prediction assumes a single row of
alternating cores exactly on the sampling line and a spacing fixed at its
far-wake value, whereas the sampling line crosses a wake that meanders laterally, and
the cores are closer together near the body than after they have convected and spread.
Both effects add reversals. What the check establishes is that the number of cores in
the near wake is consistent with the measured shedding frequency to within a small
integer factor --- enough to rule out a global oscillation with no resolved cores,
which is the alternative it exists to exclude, and not enough to serve as a
quantitative validation.

This must be read as a consistency check rather than a precise measurement, and the
reason is worth stating because it is a genuine weakness of the diagnostic. The
reversal count depends on how the wake is sampled. A naive approach that selects
cells near a line and sorts them by $x$ inflates the count severely, because it
interleaves cells at different $y$ and so crosses repeatedly between vortex cores
and the shear layers between them: on this field that approach gives anywhere from
$1$ to $73$ reversals as the selection band is widened from $0.02$ to $0.4$
diameters, which is a metric measuring its own sampling window rather than the
flow. The figures quoted above instead bin the wake into uniform $x$ intervals and
take the cell nearest the sampling line in each, giving a well-defined
one-dimensional signal. An independent reimplementation of that binned procedure,
differing in bin count and tie-breaking, returned $6$, $7$ and $9$ reversals on the
same three lines. The count is therefore reproducible to within a few reversals
rather than to the digit, and no argument here depends on more precision than that.

\paragraph{Shedding statistics.} The quantities below are computed over the second
half of the available record, $t\in[\cnsShedTStart, \cnsShedTEnd]$, so that the
initial transient is excluded. They are meaningful only if that window lies within
a saturated shedding regime; where the run did not reach saturation, the entries
are marked as unavailable rather than filled with numbers that would describe a
developing transient.

\begin{center}
\footnotesize
\begin{tabular}{lrl}
  \toprule
  Quantity & Computed & Literature ($\Reinf=200$) \\
  \midrule
  Analysis window & $t\in[\cnsShedTStart,\ \cnsShedTEnd]$ & --- \\
  Completed cycles in window & $\cnsShedCycles$ & --- \\
  Saturation gate ($\ge5$ cycles) & \textbf{\cnsShedSaturated} & --- \\
  Shedding period & $\cnsShedPeriod$ & $\approx 5.3$ \\
  Period spread ($\cnsShedEstimates$ estimates) & $\cnsShedSpread$
    ($\cnsShedSpreadPct\,\%$) & --- \\
  Strouhal number $St=fD/U_\infty$ & $\cnsShedStrouhal$ & $\approx 0.19$--$0.20$ \\
  Mean $C_D$ & $\cnsShedCdMean$ & $\approx 1.3$--$1.4$ \\
  $C_D$ peak-to-peak & $\cnsShedCdPtp$ & --- \\
  Mean $C_L$ & $\cnsShedClMean$ & $0$ by symmetry \\
  $C_L$ amplitude & $\pm\cnsShedClAmp$ & $\approx 0.6$--$0.7$ \\
  $C_L$ rms & $\cnsShedClRms$ & --- \\
  \bottomrule
\end{tabular}
\end{center}

The analysis window contains $\cnsShedCycles$ shedding cycles, and the
saturated-regime test (at least five completed cycles in the window) returns
\textbf{\cnsShedSaturated}. That qualifier governs how the table above should be
read: with fewer than a few cycles a nominal frequency estimate is an artefact of
the window length rather than a property of the flow, so the Strouhal figure is
only quotable when the test passes.

The Strouhal number is obtained from linearly interpolated zero crossings of the
lift signal, which requires no assumption about spectral resolution; with
$D=U_\infty=1$ the Strouhal number is numerically the frequency. These values come
from \texttt{tools/transient\_status.py}, the same tool used to monitor the running
case, so the statistics reported here are the ones the run-time diagnostics reported.

A second tool, an FFT-based estimator, was used earlier in the run to produce the
periodogram of \cref{fig:cyl200-spectrum}, and it must not be read as an independent
confirmation of the table: it was last run at $t\approx11$, deep in the growth phase,
where it returned $St=0.228$ and $\overline{C_D}=1.05$. Those figures describe a
pre-saturation record and are superseded by the zero-crossing values quoted above,
which are measured on the saturated window. The spectrum figure is retained because it
shows a single dominant peak rather than a broadband signal, which is a qualitative
statement that survives the frequency shift, but the numerical values in the table
come from one authority only. The lift oscillates about zero at the shedding
frequency while the drag oscillates about its mean at twice that frequency, which
is the standard signature of alternating vortex shedding: each shed vortex
produces a lift excursion of one sign but a drag excursion regardless of sign.
This behaviour is visible directly in \cref{fig:cyl200-forces}, and the wake
structure producing it in \cref{fig:cyl200-wake}.

\begin{figure}[htbp]
  \centering
  \begin{subfigure}{0.49\textwidth}\centering
    \cnsfig{cylinder_m010_laminar_re200_forces.png}{\linewidth}
    \caption{Force histories.}
  \end{subfigure}\hfill
  \begin{subfigure}{0.49\textwidth}\centering
    \cnsfig{cylinder_m010_laminar_re200_residuals.png}{\linewidth}
    \caption{Transient residual history.}
  \end{subfigure}
  \caption{Cylinder at $\Reinf=200$: post-transient $C_L$ and $C_D$ against
  physical time from \texttt{forces.csv}, and the total transient residual per
  physical step from \texttt{residuals.csv}. The lift oscillates at the shedding
  frequency and the drag at twice that frequency.}
  \label{fig:cyl200-forces}
\end{figure}

\begin{figure}[htbp]
  \centering
  \cnsfig{cylinder_m010_laminar_re200_spectrum.png}{0.62\textwidth}
  \caption{Cylinder at $\Reinf=200$: power spectrum of the post-transient lift
  signal from \texttt{forces.csv}, with the dominant peak marked. The peak
  frequency is the shedding frequency used for the Strouhal number.}
  \label{fig:cyl200-spectrum}
\end{figure}

\begin{figure}[htbp]
  \centering
  \begin{subfigure}{0.49\textwidth}\centering
    \cnsfig{cylinder_m010_laminar_re200_vorticity.png}{\linewidth}
    \caption{Vorticity, clipped to $[-5,5]$.}
  \end{subfigure}\hfill
  \begin{subfigure}{0.49\textwidth}\centering
    \cnsfig{cylinder_m010_laminar_re200_velocity.png}{\linewidth}
    \caption{Velocity magnitude.}
  \end{subfigure}\\[6pt]
  \begin{subfigure}{0.49\textwidth}\centering
    \cnsfig{cylinder_m010_laminar_re200_pressure.png}{\linewidth}
    \caption{Pressure.}
  \end{subfigure}\hfill
  \begin{subfigure}{0.49\textwidth}\centering
    \cnsfig{cylinder_m010_laminar_re200_mach.png}{\linewidth}
    \caption{Mach number.}
  \end{subfigure}
  \caption{Cylinder at $\Reinf=200$: post-transient wake visualisations from
  \texttt{field\_final.vtu}. Vorticity is clipped to the recommended $[-5,5]$
  range, without which the near-wall vorticity extremes compress the wake into a
  single colour. The alternating pattern of opposite-signed vortices is the
  vortex street; the low-pressure cores in panel (c) mark the individual
  vortices.}
  \label{fig:cyl200-wake}
\end{figure}
